\documentclass{article}

\usepackage{amssymb,amsmath,mleftright}

\begin{document}

General model assumptions:
\begin{itemize}
\item the fluid density and viscosity are constant in time and space (implies incompressibility);
\item the fluid flow is laminar;
\item the column is operated under constant conditions (e.g., temperature, flow rate);
\item the process is isothermal (i.e., there are no thermal effects);
\item diffusion does not depend on the concentration of the components, viscosity of the fluid, or pressure;
\item the fluid in the particles is stagnant (i.e., there is no convective flow);
\item the porous particles are spherical, rigid, and do not move;
\item the particles are homogeneous and of uniform porosity;
\item all components access the same particle volume;
\item the partial molar volumes are the same in mobile and solid phase;
\item the binding model parameters do not depend on pressure and are constant along the column;
\item the solvent is not adsorbed;
\end{itemize}


Specific model assumptions:
\begin{itemize}
\item the fluid only flows in the radial direction of the column (i.e., there is no flow in the axial and angular direction);
\item the column is a hollow cylinder with inner radius $R^{\mathrm{in}}$ and outer radius $R^{\mathrm{out}}$;
\item the particles form a continuum inside the column (i.e., there is interstitial and particle volume at every point in the column);
\item the column is radially symmetric and homogeneous (i.e., concentration profiles and parameters only depend on the axial position);
\item the interstitial liquid phase concentration is spatially constant on the outer particle surface;
\item the pore and surface diffusion are infinitely fast. That is, we assume $D_{i}^{\mathrm{p}} = D_{i}^{\mathrm{s}} = \infty$;
\end{itemize}


\section*{Radial Flow Lumped Rate Model with Pores}
Consider a hollow cylindrical column with inner radius $R^{\mathrm{in}} > 0$ and outer radius $R^{\mathrm{out}} > R^{\mathrm{in}}$ packed with spherical particles, and observed over a time interval $(0, T^{\mathrm{end}})$.
In the interstitial volume, mass transfer is governed by the following convection-diffusion-reaction equations in $(0, T^\mathrm{end})\times (R^\mathrm{in}, R^\mathrm{out})$ and for all components $i\in\{1, \dots, N^{\mathrm{c}} \}$
\begin{align}
\varepsilon^{\mathrm{c}} \frac{\partial c^{\mathrm{b}}_i}{\partial t} = - \frac{v}{\rho} \frac{\partial \left( \varepsilon^{\mathrm{c}} c^{\mathrm{b}}_i \right)}{\partial \rho} + \frac{1}{\rho} \frac{\partial}{\partial \rho} \left( \rho \varepsilon^{\mathrm{c}} D^{\mathrm{rad}}_{i} \frac{\partial c^{\mathrm{b}}_i}{\partial \rho} \right)- \left(1 - \varepsilon^{\mathrm{c}} \right) \frac{3}{R^{\mathrm{p}}} k^{\mathrm{f}}_{i} \left(c^{\mathrm{b}}_i - \left. c^{\mathrm{p}}_{i} \right|_{r = R^{\mathrm{p}}} \right),
\end{align}
with boundary conditions

\begin{alignat}{2}
v c_{\mathrm{in},i} &= \left.\left( v c^{\mathrm{b}}_i - D^{\mathrm{rad}}_{i} \frac{\partial c^{\mathrm{b}}_i}{\partial \rho}\right)\right|_{\rho=R^{\mathrm{in}}} & &\qquad\text{on }(0, T^{\mathrm{end}}),\\
               0 &= - D^{\mathrm{rad}}_{i} \left. \frac{\partial c^{\mathrm{b}}_i}{\partial \rho} \right|_{\rho=R^{\mathrm{out}}} & &\qquad\text{on }(0, T^{\mathrm{end}}).
\end{alignat}
Here, $t\in (0, T^{\mathrm{end}})$ is the time coordinate, $\rho\in (R^{\mathrm{in}}, R^{\mathrm{out}})$ is the radial coordinate, $R^\mathrm{p}> 0$ is the particle radius, $c^{\mathrm{b}}_i\colon (0, T^\mathrm{end})\times (R^\mathrm{in}, R^\mathrm{out}) \mapsto \mathbb{R}$ is the bulk liquid concentration, $Q> 0$ is the volumetric flow rate, $\varepsilon^{\mathrm{c}}\in (0, 1)$ is the column porosity, $v:= \frac{Q}{2 \pi L}$ is the velocity coefficient, $D^\mathrm{rad}_i\geq 0$ is the radial dispersion coefficient, and $k^\mathrm{f}_{i}\geq 0$ is the film diffusion coefficient.
In the particles, mass transfer is governed by reaction equations in $(0, T^\mathrm{end}) \times (R^\mathrm{in}, R^\mathrm{out})$ and for all components

\begin{align}

          \varepsilon^{\mathrm{p}}\frac{\partial c^{\mathrm{p}}_{i}}{\partial t} 
          &=\frac{3}{R^{\mathrm{p}}} k^\mathrm{f}_{i} \left( c^{\mathrm{b}}_{i} - c^{\mathrm{p}}_{i} \right) - \left( 1 - \varepsilon^{\mathrm{p}} \right) f^{\mathrm{bind}}_{i}\left( \vec{c}^{\mathrm{p}}, \vec{c}^{\mathrm{s}} \right),\\

          \frac{\partial c^{\mathrm{s}}_{i}}{\partial t}
          &= f^{\mathrm{bind}}_{i}\left( \vec{c}^{\mathrm{p}}, \vec{c}^{\mathrm{s}} \right).
\end{align}

Here, $c^{\mathrm{p}}_{i}\colon (0, T^\mathrm{end}) \times (R^\mathrm{in}, R^\mathrm{out}) \mapsto \mathbb{R}$ is the particle liquid concentration, $c^{\mathrm{s}}_{i}\colon (0, T^\mathrm{end}) \times (R^\mathrm{in}, R^\mathrm{out}) \mapsto \mathbb{R}$ is the particle solid concentration, $f^\mathrm{bind}_{i}$ is the adsorption isotherm function, $\vec{c}^\mathrm{p}$ is the particle liquid components vector, and $\vec{c}^\mathrm{s}$ is the particle solid components vector.
Consistent initial values for all solution variables (concentrations) are defined at $t = 0$.
\end{document}